% #############################################################################
% This is Chapter 2
% !TEX root = ../main.tex
% #############################################################################
% Change the Name of the Chapter i the following line
\fancychapter{Fundamental Concepts}
% The following line allows to ref this chapter
\label{chap:back}

This chapter introduces the linguistic and technical background
needed to understand the modeling choices made in the remainder of the thesis.
It first describes the main differences between European and Brazilian
Portuguese (\cref{sec:epbp-differences}) and then reviews the language-modeling
ideas on which modern NLP systems are built (\cref{sec:language-models}). The
chapter then presents the Transformer architecture (\cref{sec:transformers}),
the main language-model families that can be adapted through fine-tuning
(\cref{sec:model-families}), and the task formulations used for Portuguese
variety processing (\cref{sec:task-formulations}). Finally, it introduces the
adaptation techniques needed later in the thesis, from supervised fine-tuning to
reinforcement-learning-based optimization (\cref{sec:fine-tuning,sec:rl-reward-ft}).

\section{Differences Between European Portuguese and Brazilian Portuguese}
\label{sec:epbp-differences}

Publicly available Portuguese text is strongly skewed toward the Brazilian variant, which means
that models pretrained on Web-scale corpora tend to absorb BP lexical,
orthographic, and morphosyntactic preferences more readily than EP ones. This
imbalance matters in two ways. First, it makes automatic variety
identification harder in borderline cases where the distinction depends on
subtle cues rather than on unmistakable lexical markers. Second, it increases
the risk that rewriting systems default to BP-leaning outputs even when EP is
explicitly requested [Enhancing Portuguese Variety Identification with
Cross-Domain Approaches; From Brazilian Portuguese to European Portuguese]. For
this reason, an adaptation corpus must expose
contrasts that are visible in written text, yet nuanced enough that a generic
Portuguese model cannot be assumed to handle them reliably without targeted
fine-tuning [A Neural Approach to Language Variety Translation; BP2EP -
Adaptation of Brazilian Portuguese Texts to European Portuguese].

The contrast between EP and BP is not limited to isolated vocabulary
substitutions. Prior literature on paraphrastic correspondences between the two varieties
shows that many differences appear as recurrent alternatives in how equivalent
meanings are expressed, including progressive constructions, clitic placement,
pronoun realization, determiner use, forms of address, and multiword lexical
choices \citep{barreiro-mota-2018-paraphrastic}.
\citet{Cyrino_2001} also discusses contrasts involving null objects and verb phrase
ellipsis, relating them to changes in the verbal structure of BP sentences.

The differences below illustrate several phenomena
that are especially relevant for variety identification and controlled
rewriting:

\begin{itemize}
\item \textbf{Orthography and morphology.} Differences in spelling persist,
although many contrasts were reduced by the Portuguese Orthographic Agreement.
In particular, older EP texts may still retain muted consonants such as
\textit{projecto}, \textit{acto}, or \textit{baptismo}, whereas contemporary
standardized spelling typically drops them, yielding \textit{projeto},
\textit{ato}, and \textit{batismo}. Such forms can still appear in legacy
corpora and archival material, but are less reliable markers in
post-Agreement data \citep{costa-jussa-etal-2018-neural}.

\item \textbf{Progressive and aspectual periphrases.} EP typically favors the
progressive with \textit{estar a} + infinitive (e.g.,\textit{Ele estava a
correr}), whereas BP strongly prefers \textit{estar} + gerund (e.g., \textit{Ele estava correndo}) \citep{barreiro-mota-2018-paraphrastic}.

\item \textbf{Clitic placement and word order.} EP commonly places clitics
postverbally with hyphenation in affirmative contexts (e.g., \textit{Ele
viu-me na rua}), while BP more often favors preverbal placement without
hyphenation (e.g., \textit{Ele me viu na rua}) \citep{barreiro-mota-2018-paraphrastic}.

\item \textbf{Subject pronoun realization (pronoun drop).} Both varieties
allow null subjects, but BP more frequently includes subject pronouns in
contexts where EP more readily omits them. For example, EP often allows a null
subject in a sentence such as \textit{Fui ao supermercado ontem}, whereas BP
more often uses an overt pronoun (e.g., \textit{Eu fui ao supermercado ontem})
\citep{barreiro-mota-2018-paraphrastic, costa-jussa-etal-2018-neural}.

\item \textbf{Pronoun preferences.} BP frequently uses \textit{você} as the
default second-person address form, whereas EP more often uses \textit{tu} in
informal contexts, or omits the overt subject altogether. In practice, this
creates frequent rewrite pairs between EP second-person verb morphology and BP
third-person morphology with \textit{você} \citep{barreiro-mota-2018-paraphrastic, costa-jussa-etal-2018-neural}.

\item \textbf{Determiners with possessives and zero determiners.} In BP,
definite articles before possessives are often omitted (e.g., \textit{Vendi
meu carro}), whereas EP more typically requires them (e.g., \textit{Vendi o meu
carro}). More broadly, BP omits more determiners in contexts where EP prefers
an overt article \citep{barreiro-mota-2018-paraphrastic, Cyrino_2001}.

\item \textbf{Lexical and multiword preferences.} Some of the most robust
distinctions in written corpora arise from lexical choice and multiword
expressions, such as \textit{comboio}/\textit{trem},
\textit{fato}/\textit{terno}, or \textit{toda a gente}/\textit{todo o mundo}
\citep{barreiro-mota-2018-paraphrastic, costa-jussa-etal-2018-neural}. Apart from these more familiar
alternations, lexical variation also includes strong variety preferences and
false cognates. For instance, \textit{nomeadamente} is very frequent in EP and
rare in BP, while \textit{propina} most often means \textit{fee} in EP but is
more commonly interpreted as \textit{bribe} in BP \citep{costa-jussa-etal-2018-neural}.
\end{itemize}

From an NLP perspective, some of these differences provide strong surface cues
for variety identification, especially lexical preferences and a subset of
orthographic contrasts. Others are more structural and are therefore
particularly important for rewriting, because the model must alter sentence
organization rather than merely substitute isolated words.

These differences are consistent with the quantitative analysis reported in
\citep{alves-2025-information}, which studies EP/BP variation in the FRMT dataset using information-theoretic
measures. The main finding is that lexical choice accounts for the largest
divergence between the two varieties, but the study also identifies systematic
grammatical preferences, particularly in subject-pronoun use, progressive
constructions, and clitic-related patterns. Additional analyses based on
part-of-speech distributions, dependency relations, and word-order entropy
suggest that EP and BP share most of their syntactic backbone, while still
diverging in recurrent local preferences that are salient enough to affect
automatic processing.

To make these contrasts visible, \cref{tab:frmt-ep-bp-examples} presents a small
set of representative FRMT pairs. The examples are illustrative and were selected to mirror the divergences highlighted by \citet{alves-2025-information}.

\begin{table}[t]
\centering
\small
\setlength{\tabcolsep}{6pt}
\begin{tabular}{>{\raggedright\arraybackslash}p{0.20\linewidth} >{\raggedright\arraybackslash}p{0.35\linewidth} >{\raggedright\arraybackslash}p{0.35\linewidth}}
\toprule
\textbf{Phenomenon} & \textbf{EP example} & \textbf{BP example} \\
\midrule
Lexical choice &
\textit{O autocarro está atrasado por causa do trânsito.} &
\textit{O ônibus está atrasado por causa do trânsito.} \\
Progressive aspect &
\textit{... ficámos a conversar durante algum tempo.} &
\textit{... ficamos conversando por um tempo.} \\
Forms of address and object realization &
\textit{Se fores à mercearia, traz-me um pacote de bolachas e um sumo de laranja...} &
\textit{Se você for ao mercado, traz um pacote de biscoitos e um suco de laranja para mim...} \\
Possessive determiner and progressive aspect &
\textit{A minha irmã está a planear um jantar romântico...} &
\textit{Minha irmã está planejando um jantar romântico...} \\
\bottomrule
\end{tabular}
\caption[Representative FRMT examples for EP/BP contrasts]{Representative
sentence pairs from FRMT illustrating the lexical and morphosyntactic tendencies discussed
above.}
\label{tab:frmt-ep-bp-examples}
\end{table}

Taken together, these contrasts show why EP/BP processing is not well modeled
as a purely lexical substitution problem and motivate the need for variety-aware
adaptation. The next subsection therefore introduces language
models as the main computational framework on top of which such adaptation is
built.

\section{Language Models}
\label{sec:language-models}

Language modeling is the task of assigning probabilities to sequences of text.
Given a sequence of tokens $x_1, x_2, \ldots, x_n$, a language model estimates
the probability of the full sequence by factorizing it into conditional
probabilities, typically of the form
$p(x_1, \ldots, x_n) = \prod_t p(x_t \mid x_{<t})$. In practice, this means
that the model learns which tokens are likely to occur in a given context. In
autoregressive settings, this objective appears directly as next-token
prediction, whereas in other self-supervised settings the model may be trained
to reconstruct masked or corrupted portions of text from the surrounding
context \citep{zhao2025surveylargelanguagemodels}.

Early language models were predominantly statistical. One of the most common
approaches was the $n$-gram model, which estimates the probability of a token
from the previous $n-1$ tokens by counting how often short sequences occur in a
corpus. These models were simple and useful for capturing local regularities,
but they suffered from data sparsity and were inherently limited in their
ability to model long-range dependencies, since the context window was fixed in
advance \citep{880083}. As a result, statistical
language modeling depended heavily on smoothing techniques and large amounts of
carefully curated training data \citep{Jelinek1997StatisticalMF}.

Neural language models \citep{10.5555/944919.944966} addressed part of this problem by replacing explicit
count-based statistics with learned distributed representations. Instead of
treating each context as an isolated symbolic pattern, models used feed-forward networks to learn dense vector
embeddings for tokens and use trainable nonlinear functions to generalize
across similar contexts. Recurrent neural networks, and later gated variants
such as Long-Short Term Memory (LSTM) networks, further improved the ability to handle variable-length context
and longer dependencies, although training them efficiently at scale remained a
challenge \citep{Mikolov2010RecurrentNN}.

The current paradigm is strongly shaped by transfer learning, in which a model is first pretrained on large unlabeled corpora and then adapted to downstream tasks. This shift closely tied to the emergence of Transformer architectures, which use self-attention instead of recurrent processing and
made long-range context more directly accessible, while also enabling more
efficient large-scale training \citep{46201}. Pre-training on
large unlabeled corpora followed by task-specific adaptation then became the
standard recipe for deploying language models across a wide range of NLP tasks.

BERT \citep{devlin2019bertpretrainingdeepbidirectional} helped consolidate Transformer-based encoder models for NLP by showing
that bidirectional masked-token pretraining could transfer effectively to
understanding-oriented tasks such as classification and retrieval with minimal
task-specific changes. As parameter counts, data diversity, and context
windows continued to grow, decoder-oriented models also became increasingly
powerful generators. At large scales, these models began to
display strong prompt-based generalization, in-context learning, and other
emergent capabilities \citep{10.5555/3495724.3495883, zhao2025surveylargelanguagemodels}. 

In parallel, encoder-decoder models helped reinforce a unified
sequence-to-sequence view of NLP, in which many tasks can be formulated as
text-conditioned text generation rather than as entirely separate modeling
problems \citep{Raffel2019ExploringTL}.

An important design choice in modern neural language modeling concerns tokenization. Rather than operating directly on raw words or characters, many systems represent text as sequences of subword units. 
Frequent words may remain intact, whereas rarer or morphologically complex forms may be split into smaller pieces. This reduces out-of-vocabulary problems, allows lexical material to be reused across related forms, and helps models handle spelling variation and productive morphology \citep{sennrich-etal-2016-neural}.
For EP/BP processing, these relevant distinctions may appear either as full lexical substitutions or as smaller orthographic and morphological variations.

As stated before, pretrained language models are a natural starting point for
the tasks studied in this work, as they already encode a broad prior over
Portuguese usage and therefore provide a much stronger basis than task-specific
training from scratch. At the same time, they can be further adapted so that
the generic knowledge acquired during pretraining becomes more sensitive to the
distinctions required for variety identification and controlled rewriting. The
next subsection therefore turns to the Transformer architecture, which underlies
most modern pretrained language models.

\section{Transformers: Self-Attention, Encoders and Decoders}
\label{sec:transformers}

To explain in more detail what a Transformer is, it is useful to begin with the
self-attention operation. Self-attention allows each token in a sequence to mix
information from the others by computing attention weights that reflect learned
similarities between token representations. Each token is first mapped to a
continuous embedding. Since self-attention is permutation-invariant, token
order must be injected separately through positional encodings. In the
original Transformer, these are fixed sinusoidal vectors added to token
embeddings, yielding input representations
$X \in \mathbb{R}^{n \times d_{\text{model}}}$, where $n$ is the sequence
length and $d_{\text{model}}$ is the hidden dimension. The value of $n$ is in
practice bounded by the model's context window, that is, the maximum number of
tokens it can attend to in one forward pass \citep{46201}.

Given token representations $X$, which already include positional information,
queries, keys, and values are obtained through learned linear projections,
\[
Q = XW^Q, \qquad K = XW^K, \qquad V = XW^V.
\]
The scaled dot-product attention used throughout the architecture is then given
by
\[
\mathrm{Attention}(Q, K, V) =
\mathrm{softmax}\left(\frac{QK^\top}{\sqrt{d_k}}\right)V.
\]
This operation replaces each token representation by a similarity-weighted
combination of the others. Multi-head attention repeats this computation with
different learned projections and concatenates the resulting representations,
allowing the model to capture several interaction patterns in parallel
\citep{46201}.

Although many later architectures are based on self-attention, the original
Transformer was proposed as an encoder-decoder model for sequence-to-sequence
tasks. An encoder is a stack of Transformer blocks, each formed by a
self-attention sublayer followed by a position-wise feed-forward network, with
residual connections and layer normalization supporting stable training. The
feed-forward network is position-wise in the sense that the same multilayer
perceptron is applied independently to each token representation. Since encoder
self-attention is bidirectional, every position can attend to every other
position in the source sequence, producing contextual representations that are
useful both for text understanding and for conditioning downstream generation
\citep{46201}.

A decoder is also a stack of Transformer blocks, but it introduces two
mechanisms required for generation. First, masked self-attention applies a
causal mask that prevents each position from attending to future tokens, so the
model generates from left to right. Second, a cross-attention sublayer allows
the decoder to attend to the encoder outputs, making target generation
explicitly conditioned on the source sequence \citep{46201}.

This architectural template gives rise to distinct model families, which differ in training objectives
and in how naturally they support classification and text generation. The next subsection reviews these
families and motivates the model selection for the initial experiment performed in this work.

\section{Language Model Families Amenable to Fine-Tuning}
\label{sec:model-families}

Modern pretrained language models rely on the Transformer architecture, but
they differ in how the architecture is used and in the pretraining objectives
that shape their representations. The main families relevant to
downstream adaptation are encoder-only, decoder-only, and encoder-decoder
models.

Encoder-only models, such as BERT, are bidirectional encoders trained to
recover masked tokens or spans from their surrounding context
\citep{devlin2019bertpretrainingdeepbidirectional}. Since each position can
attend to both left and right context, these models are especially effective at
building contextual representations for tasks such as classification, retrieval,
or span prediction. In practice, they are often adapted by attaching a
task-specific prediction head on top of the encoder. However, they are not
generative by default, which makes them less convenient
when one wants a single model family to cover both understanding-oriented and
generation-oriented tasks.

Decoder-only models use a single Transformer stack with causal masking and
define an autoregressive distribution over continuations
\citep{10.5555/3495724.3495883}. As their native operation is text generation, given a prefix sequence the model defines the probability of each next token conditioned on the tokens that precede it. 
This makes decoder-only models a natural fit for open-ended
generation. When they undergo additional adaptation through instruction tuning,
they can also become well suited to task following and prompting-based interaction.
These models can be used for classification as well, either by generating a label directly
or by scoring candidate labels according to their conditional likelihood. 

Encoder-decoder models combine a bidirectional encoder with an autoregressive
decoder. The encoder constructs contextual representations of the source
sequence, and the decoder generates the target sequence while attending to that
source representation. This formulation is the classical choice for machine
translation and other conditional generation tasks, and it also supports the
text-to-text view in which multiple NLP problems are expressed as mappings from
an input sequence to an output sequence \citep{Raffel2019ExploringTL}. 

All three families operate over tokenized text within a finite context window
and can benefit from transfer learning through downstream fine-tuning. The key
question is therefore not whether these families can in principle represent the
tasks studied, but how naturally their pretraining setup aligns with the
desired use case. Encoder-only models are particularly strong for
understanding-oriented prediction and decoder-only models provide a flexible
generative interface. As this work focuses on both bidirectional rewriting and
variety-aware prediction, encoder-decoder models provide a particularly
coherent starting point for the modeling choices introduced in the following
sections.

Within this encoder-decoder space, we considered a model from the T5Gemma 2 family, most notably the
\texttt{t5gemma-2-4b-4b} checkpoint. T5Gemma 2 is presented as a family of
encoder-decoder models obtained by adapting pretrained decoder-only Gemma
models and then continuing pretraining with a UL2-style objective, which mixes
several denoising objectives, including span corruption and prefix-style
prediction, so that the model is exposed both to bidirectional understanding
and to autoregressive generation 
\citep{zhang2025t5gemma2seeingreading, 52876}.

This matches the main setups used in the present work reasonably well. The
dominant task is source-conditioned rewriting, where the model receives a BP or
EP sentence and must generate its counterpart in the other variety. In
addition, some of the setups used encode classification-like
information directly in the decoder target, for instance by making the decoder
begin with a short variety label before producing the rewritten sentence. In
that sense, the appeal of T5Gemma 2 in this research is not mainly its
multimodal or very long-context capability, which are not central to our
experiments, but rather its compatibility with conditional generation and
structured text targets.

T5Gemma 2 also introduces architectural choices that are worth noting. For example, these models share word embeddings
across the encoder and the decoder and replace the usual separation between decoder
self-attention and cross-attention with a merged attention module
\citep{zhang2025t5gemma2seeingreading}. The first choice reduces parameter
duplication, while the second aims to improve efficiency without removing the
decoder's access to the encoded source. Recent work on adapting Gemma models
from decoder-only to encoder-decoder form further argues that this architecture
offers a favorable trade-off between richer source representations, efficient
inference, and strong downstream fine-tuning behavior \citep{Zhang2025EncoderDecoderGI}.

The next step is to make explicit how tasks are represented for such a model. The
following section therefore defines variety identification and variety
rewriting as task formulations.

\section{Task Formulations for Portuguese Variety Processing}
\label{sec:task-formulations}

In this work, Portuguese variety processing is organized around two
complementary formulations. One asks the model to decide which variety
a text belongs to, while the other asks it to produce a version of a sentence in
a requested variety. This distinction is useful as both tasks depend on
the same linguistic contrast, but expose that contrast through different output
structures.

The standard formulation of variety identification is supervised
classification where, given a text segment, the system selects one label from a
small set of possible varieties, such as EP or BP
\citep{article, sousa2025enhancingportuguesevarietyidentification}. However, in a generative
model, the same task can be represented as generation over a restricted
output space, where the decoder is expected to produce only a short variety
label. This does not change the underlying task, but changes the model's
interface. In a unified sequence-to-sequence setting, this can be implemented
with label-only targets for identification examples, or with structured targets
where the first generated token is a variety label. If the target contains both
a label and a longer sentence, the learning signal can also be restricted to
the label positions for identification-specific objectives. 

On the other hand, variety rewriting is naturally formulated as conditional generation. The model
receives a source sentence together with a target-variety
specification and generates an equivalent sentence in the requested variety. At
the task-formulation level, EP/BP rewriting is close to language-variety
translation, where the source and target are not different languages but
national or regional varieties of the same language
\citep{costa-jussa-etal-2018-neural,sanches2024brazilianportugueseeuropeanportuguese}.
The rewriting formulation is also related to paraphrastic reformulation, since
many valid conversions between EP and BP involve meaning-preserving
alternatives at the lexical, phrasal, or syntactic level
\citep{barreiro-mota-2018-paraphrastic}.

A text-to-text formulation provides a common representation for these two
tasks. In this setting, both inputs and outputs are token sequences, and the
difference between tasks lies mainly in the expected target sequence. For
identification, the target is a short label and, for rewriting, the target is a
full sentence. This type of unified formulation has been used to express many
NLP problems within a single sequence-to-sequence framework
\citep{Raffel2019ExploringTL}. For the present work, the motivation is not
only interface simplicity. The two tasks rely on overlapping variety-sensitive
evidence, since the cues that support an EP/BP decision are often the same cues
that must be preserved or modified during rewriting. Sharing model parameters
across tasks can therefore allow identification data to reinforce what signals
variety, while rewriting data teaches how those signals should change when a
target variety is requested \citep{10.1023/A:1007379606734}. This does not imply that a
multitask model must outperform a dedicated single-task model on every metric,
since single-task models can allocate all capacity to one objective. However,
one of the benefits is to obtain a single model that can perform both
variety-aware prediction and controlled rewriting.

The desired variety can also be made explicit through conditioning. In
region-aware translation benchmarks, the target region or variety is part of
the task specification, so the model is expected to generate an output that is
appropriate not only in meaning but also in regional form
\citep{riley2023frmtbenchmarkfewshotregionaware}. In text-to-text models, this
kind of conditioning can be represented in several ways. For example, a task
prefix can be placed in the input to indicate the operation to perform. A
control token can specify the desired attribute of the output, such as the
target variety. In addition, a structured target can require the decoder to
emit information in a fixed order, for example a variety label
followed by the generated sentence \citep{Raffel2019ExploringTL}. 
The exact labels, token choices, and dataset conversion rules are
implementation details and are therefore described later in the methodology
chapter.

\section{Fine-Tuning Pretrained Language Models}
\label{sec:fine-tuning}

The previous sections described how Portuguese variety processing can be
expressed as text-to-text prediction and why encoder-decoder models are a
natural backbone for this setting. However, pretraining alone does not guarantee
that a model will follow the specific behavior required by a downstream task.
Even when a pretrained model has broad linguistic knowledge, it may not
consistently distinguish EP from BP, preserve meaning during rewriting, or obey
the requested target variety. Fine-tuning is the adaptation step that addresses
this gap by continuing training on data that is closer to the intended use case.

In supervised fine-tuning, the model is trained on input-output examples and
the parameters are updated to increase the likelihood of the desired targets.
For a text-to-text model, this means that both classification-style examples and
rewriting examples can be represented as sequence prediction problems
\citep{Raffel2019ExploringTL}. A related form of adaptation is continued
pretraining or domain-adaptive pretraining, where the model is exposed to
additional in-domain text while keeping a self-supervised objective rather than
a task-specific target [Don't Stop Pretraining: Adapt Language Models to
Domains and Tasks]. The difference is that supervised fine-tuning teaches the
model the desired mapping, whereas continued pretraining mainly adapts the
model's internal representations to a domain or variety of text.

A central design choice is how many parameters should be updated during
adaptation. Full fine-tuning updates the entire pretrained model. This gives the
optimizer maximum flexibility and can be useful when the target task or domain
differs substantially from pretraining, but it is expensive in memory and
computation, and may overfit when the supervision is narrow. Parameter-Efficient
Fine-Tuning (PEFT) instead freezes most or all of the pretrained backbone and
learns a much smaller set of task-specific parameters. Adapter-based methods
insert small trainable modules into the network while keeping the original
model fixed, allowing task adaptation with far fewer updated parameters than
full fine-tuning \citep{houlsby2019parameterefficienttransferlearningnlp}.

LoRA follows the same parameter-efficient motivation, but implements it by
learning low-rank updates to selected weight matrices rather than by inserting
separate adapter layers. If a pretrained weight matrix is denoted by
\(W \in \mathbb{R}^{d \times k}\), LoRA keeps \(W\) frozen and learns an update
\(\Delta W = BA\), where \(B \in \mathbb{R}^{d \times r}\),
\(A \in \mathbb{R}^{r \times k}\), and \(r\) is much smaller than \(d\) and
\(k\). The effective adapted weight is therefore \(W + \Delta W\), but only the
low-rank factors are trained \citep{hu2021loralowrankadaptationlarge}. This
substantially reduces the number of trainable parameters and makes it possible
to adapt larger models under limited hardware. Quantized variants extend this
idea by storing the frozen backbone in low precision while training the adapter
parameters in higher precision [QLoRA: Efficient Finetuning of Quantized LLMs].

Fine-tuning also need not occur in a single step. In staged or sequential
fine-tuning, adaptation is divided into phases that may use different datasets,
objectives, or levels of supervision. This can be useful when one phase is meant
to teach broad task behavior, another to specialize the model to cleaner or more
specific data, and a later phase to optimize preferences or reward-based
criteria. The advantage is that each phase can emphasize a different learning
signal; the risk is that later phases can overwrite behavior learned earlier,
especially when the data distribution or objective changes substantially
\citep{pfeiffer-etal-2021-adapterfusion}. In this thesis, the concept is
important because variety identification, bidirectional rewriting, and
reward-guided refinement are related but not identical forms of adaptation. The
next section therefore introduces the background needed when the final
adaptation signal is expressed as a reward rather than as a fixed reference
target.

\section{Reinforcement Learning and Reward-Based Fine-Tuning}
\label{sec:rl-reward-ft}

% Scope for this subsection:
% - Start from the general idea of optimizing model behavior with rewards or
%   preferences rather than only token-level supervision.
% - Cover preference-based fine-tuning and RL-style optimization at a level
%   broad enough to support the later Stage C discussion.
% - Include the idea of combining multiple reward components, since this is
%   central to the thesis setup.
% - Explain reward weighting, reward shaping, and the possibility of mixing
%   supervised and RL losses only conceptually.
% - Do not put the thesis-specific reward design here (BLEU/TER/first-token
%   weights, copy penalty values, exact candidate-generation settings); those
%   belong in methods and experiments.
